\section{Introduction}
\label{sec:introduction}

The field of AI creative systems rests on an implicit assumption: domain expertise is
the primary quality variable. Add a puzzle-designer persona, get better puzzle design.
Add a legal-analyst persona, get better legal analysis. Add a marketing-expert persona,
get better marketing copy. This assumption shapes how practitioners write prompts, how
researchers parameterize personas, and how creative AI systems are evaluated. It is
wrong --- wrong not in degree but in kind, in a specific and actionable way.

Three empirical studies across authoring and reviewing roles in puzzle-hunt design
reveal a consistent three-tier structure in which domain expertise occupies only the
outermost and most substitutable tier. The irreplaceable tier is not expertise at all.
It is an orientation: a stance toward the person on the other end of the creative
interaction --- the solver who will encounter the puzzle, the reader who will consult the
review, the user who will inhabit the artifact. This orientation, expressible as a
single sentence, outperforms any domain label as a predictor of creative quality. Our
motivating result: the bare two-word persona ``you are an artist'' outperformed ``you
are a puzzle designer'' by 15.6 points on a 45-point weighted rubric. Game designer
scored 32.3 out of 45, programmer 24.0, doctor 17.7. The artist was the only condition
to escape mechanism gravity --- producing a structurally novel puzzle mechanism rather
than reusing the acrostic-and-indexing family that dominated all domain-labeled
conditions.

This result sets the explanatory challenge for the paper. Why does a domain-naive
label outperform a domain-specific one? The answer is not that the artist label is
superior by accident. It is that the two labels activate different tiers of the
creative quality hierarchy. ``You are a puzzle designer'' activates Tier 3 content:
domain conventions, format expectations, the surface grammar of the target artifact.
``You are an artist'' activates Tier 1 content: an orientation toward experience,
toward the receiver, toward the question of what the artifact will feel like to
encounter. When a system optimizes for Tier 3 without Tier 1, it produces work that
is formally correct and experientially inert.

The empirical decomposition that follows the motivating study identifies a three-tier
structure across both the authoring and reviewing roles. In authoring, we decomposed
the artist persona into six cognitive traits and measured quality under each possible
ablation. The experience-first trait (Tier 1) was irreplaceable: removing it caused a
catastrophic drop from 41.7 to 12.3 out of 45, the only condition to produce
categorical failure. Elegance drive and rigor formed the load-bearing Tier 2: removing
either dropped quality below a 35-point threshold associated with publishable puzzle
quality. The remaining traits --- surprise, visual thinking, systematic reasoning ---
were Tier 3 amplifiers: each added value but none was necessary for threshold crossing.

In reviewing, an eight-condition ablation study of a structured expert reviewer profile
revealed the same structure. The design philosophy and worldview component (Tier 1) was
irreplaceable: the lens-only condition (Tier 3 without Tier 1) scored 21.7, falling
below the bare baseline condition with no persona at all (23.7). Three non-redundant
operational principles formed the load-bearing Tier 2. The review lens and credentials
were Tier 3 amplifiers. When all eleven reviewer principles were included without the
underlying philosophy, quality degraded from 25.3 to 22.3 --- principles without
orientation produce noise, not signal.

The symmetry between the authoring and reviewing findings is not coincidental. It
follows from a general principle: creative quality is defined from the receiver's
perspective, enforced through operational constraints, and enriched by domain
knowledge. This ordering holds regardless of whether the role is creating or
evaluating, because in both cases the ultimate criterion is the experience of the
downstream person.

The practical implication is the OLE formula: Orientation before Lens before
Expertise. Every AI creative prompt should begin with an orientation sentence
(``Think first about what the other person will experience''), proceed to discipline
sentences (operational constraints that enforce the orientation at decision points),
and conclude with domain content (world data, style guides, precedents) as enrichment
rather than identity. Three sentences constitute the minimum sufficient creative prompt
for AI creative tasks in domains where quality is defined by receiver experience.
We predict this prescription generalizes beyond puzzle-hunt design (cross-domain validation pending).

This paper makes four contributions. First, we present a unified three-tier model of
creative quality that applies symmetrically across authoring and reviewing roles, with
specific empirical grounding for each tier in each role. Second, we formalize the OLE
formula as a practical prescription for AI creative prompt design with quantitative
support for its ordering. Third, we demonstrate empirically that a domain-naive
orientation outperforms a domain-specific expertise label in a controlled authoring
study, and explain the mechanism via the tier model. Fourth, we argue for the
generalizability of the tier structure beyond puzzle design, identifying predicted tier
assignments for business-domain creative tasks and specifying the boundary conditions
under which the ordering might differ.

The remainder of the paper proceeds as follows. Section~\ref{sec:related} positions the
tier model relative to prior work on AI persona prompting, expertise structure, and
experience-first design. Section~\ref{sec:tiermodel} presents the model with its
empirical grounding. Section~\ref{sec:evidence} presents the cross-role evidence in
detail. Section~\ref{sec:ole} formalizes the OLE formula and develops the
generalization argument. Sections~\ref{sec:discussion} and~\ref{sec:conclusion} discuss
implications and open questions.
